% =============================================================================
% Kapitel 1: Einleitung
% -----------------------------------------------------------------------------
% Was hinein gehört (Richtwert 4--6 Seiten):
%  - Motivation: Synchrotron-Röntgenbildgebung, Phasenkontrast, warum NFH
%    (schwach absorbierende Objekte, Kohärenz an PETRA III / P05).
%  - Problem: Phase Retrieval braucht ein genaues Vorwärtsmodell; die Geometrie
%    steckt in der Fresnel-Zahl Fr. Falsches Fr -> Artefakte. Autofokus =
%    Bestimmung von Fr (bzw. z01) aus den Daten.
%  - Stand: modellbasierter Autofokus (Dora et al. 2025) ist interpretierbar,
%    aber langsam (Optimierung mit vielen Rekonstruktionen).
%  - Ziel dieser Arbeit: lernbasierte Schätzung von Fr/z01 direkt aus dem
%    Hologramm; Trainingsdaten aus HoloWizard/HoloForge; optional
%    Unsicherheiten via SBI.
%  - Forschungsfragen (nummeriert, später in Kapitel 6/7 wieder aufgreifen):
%      F1 Wie genau kann ein CNN Fr/z01 aus einem einzelnen Hologramm schätzen?
%      F2 Wie verhält sich das im Vergleich zum modellbasierten Autofokus
%         (Genauigkeit vs. Laufzeit)?
%      F3 Übertragbarkeit Simulation -> Messdaten (Sim-to-Real-Gap)?
%      F4 Liefert SBI/NPE brauchbare Unsicherheiten?
%  - Beiträge der Arbeit (Stichpunkte).
%  - Aufbau der Arbeit (ein Absatz, mit \cref auf Kapitel).
% =============================================================================
\chapter{\lang{Einleitung}{Introduction}}
\label{ch:einleitung}

\section{\lang{Motivation}{Motivation}}
\label{sec:einleitung:motivation}

% Beispiel: Zitat + Abkürzung beim ersten Auftreten
\lang{%
Die \ac{nfh} nutzt die hohe Kohärenz von Synchrotronstrahlung, um auch schwach
absorbierende Objekte über ihren Phasenkontrast abzubilden. Das zugrunde liegende
Rekonstruktionsproblem, das \ac{pr}, ist nichtlinear und schlecht gestellt
\cite{fienup1982phase}. Moderne Verfahren wie die \ac{asrm} rekonstruieren das
komplexe Objekt ohne räumliche Trägerbeschränkung \cite{dora2024asrm}, setzen
aber ein genaues Modell der Aufnahmegeometrie voraus.
}{%
\Ac{nfh} exploits the high coherence of synchrotron radiation to image weakly
absorbing objects via their phase contrast. The underlying reconstruction problem,
\ac{pr}, is non-linear and ill-posed \cite{fienup1982phase}. Modern methods such
as \ac{asrm} reconstruct the complex object without a spatial support constraint
\cite{dora2024asrm}, but require an accurate model of the acquisition geometry.
}

\todo{\lang{Motivation ausformulieren.}{Elaborate the motivation.}}

\section{\lang{Problemstellung und Zielsetzung}{Problem Statement and Objectives}}
\label{sec:einleitung:ziel}

\lang{%
Der aktuelle Stand der Technik bestimmt die Fresnel-Zahl $\Fr$ durch
modellbasierte Optimierung \cite{dora2025autofocus}. Ziel dieser Arbeit ist es,
$\Fr$ bzw. den Abstand $\zOne$ mit einem neuronalen Netz direkt aus dem Hologramm
zu schätzen, um die Autofokussierung um Größenordnungen zu beschleunigen.
}{%
The current state of the art determines the Fresnel number $\Fr$ by model-based
optimisation \cite{dora2025autofocus}. This thesis aims to estimate $\Fr$ or the
distance $\zOne$ directly from the hologram with a neural network in order to
speed up autofocusing by orders of magnitude.
}

\todo{\lang{Forschungsfragen F1--F4 formulieren.}{State research questions F1--F4.}}

\section{\lang{Aufbau der Arbeit}{Outline}}
\label{sec:einleitung:aufbau}

\lang{%
\Cref{ch:grundlagen} führt die physikalischen und methodischen Grundlagen ein,
\cref{ch:stand} ordnet verwandte Arbeiten ein. \Cref{ch:methoden} beschreibt
Datenerzeugung, Modelle und Trainingsaufbau, \cref{ch:experimente} die
Experimente und Ergebnisse. \Cref{ch:diskussion} diskutiert die Befunde,
\cref{ch:fazit} fasst zusammen und gibt einen Ausblick.
}{%
\Cref{ch:grundlagen} introduces the physical and methodological background,
\cref{ch:stand} reviews related work. \Cref{ch:methoden} describes data
generation, models and the training setup, \cref{ch:experimente} the experiments
and results. \Cref{ch:diskussion} discusses the findings and \cref{ch:fazit}
concludes with an outlook.
}
